% Stage 8 restored fragment — 2.3.1 Scalar, vector, and dyadic Green functions (archive lines 3515-3952)
\subsection{Scalar, vector, and dyadic Green functions}
\label{sec:ch231}
Physically, $G$ and $\overline{\mathbf{G}}_{E}$ are the medium's response functions: each
records how a unit impulse at $\mathbf{r}'$ launches scalar or vector fields to $\mathbf{r}$ under
the same boundary completion as $\mathcal{L}_\omega$ \cite{stratton1941}. Every compact-range campaign begins with a dipole elementary check: a electrically small feed
must reproduce the transverse far-zone law before multipole catalogs are trusted on larger
apertures \cite{jackson1999}. That observation anchors the Green-function program
developed here---first the scalar monopole response, then vector potentials, then the
electric dyad that maps $\Jimp$ to $\E$ directly \cite{stratton1941}. The dyadic Green
function is the electromagnetic impulse response of the governing Maxwell
operator: the field at an observation point produced by an elemental current located
elsewhere. In uniform media it is also the continuous limit of one radiating cell in a
phased array and the local kernel that probe correction applies when near-field data are
translated to a reference surface \cite{stratton1941}. Admissible kernels are
retarded---field response follows source excitation only after finite propagation delay---so
localized impressed currents remain consistent with Maxwell causality
\cite{sommerfeld1949}. Scalar Helmholtz functions describe monopole-like propagation;
polarization coupling on vector feeds requires the full dyadic \cite{collin2001}.

The \emph{kernel formulas} developed below must be distinguished from
and the \emph{domain declarations} that make those formulas meaningful on a declared
boundary-value problem \cite{stratton1941}. Unless noted otherwise, \(\omega>0\) is fixed,
\(\Jimp\in\mathcal{J}_\Omega\) from Eq.~\eqref{eq:ch2.source-space-refined}, charge continuity
Eq.~\eqref{eq:ch1.charge-continuity-harmonic} holds, and \(\mathcal{B}_\Gamma\) is declared on
\(\Gamma=\partial\Omega\). In uniform vacuum regions the wavenumber is \(k_0\) from
Eq.~\eqref{eq:ch2.helmholtz-symbol}. Kernels are stated first in operator form on
\(\mathcal{D}(\mathcal{W}_{E,\omega})\subset\HCurl(\Omega)\) from
Eq.~\eqref{eq:ch2.electric-helmholtz-domain}; piecewise \(k(\mathbf{r})\) extensions use the
same template with patchwise coefficients \cite{stratton1941}. The harmonic solution operator
$\mathcal{G}_{E,\omega}$ introduced in Subsection~\ref{sec:ch222} is realized below by
convolution with $G$ and $\overline{\mathbf{G}}_{E}$; observation coordinate $\mathbf{r}$ and
source coordinate $\mathbf{r}'$ must remain distinct throughout every range transformation
\cite{harrington2001}.
Imagine exciting a uniform medium with the smallest source one can conceive: a unit
impulse concentrated at a single point $\mathbf{r}'$. The field that appears at every
other location $\mathbf{r}$ is the scalar Green function $G(\mathbf{r},\mathbf{r}')$---the
impulse response of the Helmholtz operator at fixed frequency
\cite{stratton1941}. Laboratory engineers never place a mathematical
$\delta$-function on a bench, yet every extended feed, patch, or volume current is assembled
from such elemental responses by superposition: each infinitesimal slice of $\Jimp(\mathbf{r}')$
radiates as if it were a point source scaled by the local source strength, and the total
field is the sum of those contributions weighted by propagation from $\mathbf{r}'$ to
$\mathbf{r}$ \cite{harrington2001}. Green functions are therefore the natural language of
propagation: they replace repeated boundary-value solves with a single propagator integral
that method-of-moments codes, near-field scanners, and imaging algorithms all discretize
\cite{stratton1941}.
The physical significance of Eq.~\eqref{eq:ch2.scalar-helmholtz-green-pde} is that it isolates
the response of the medium to the smallest source that can be imagined---a unit point excitation.
Once this elementary response is known, more complicated sources may be synthesized by
superposition \cite{jackson1999}. The scalar Green function $G(\mathbf{r},\mathbf{r}')$
inverts the Helmholtz operator $(\nabla^2+k_0^2)$ on scalar potentials; $\mathbf{r}$ denotes
the observation point and $\mathbf{r}'$ the source point, with $R=|\mathbf{r}-\mathbf{r}'|$
the geometric separation in metres. On a compact range, the scalar kernel is the monopole
reference against which probe calibration and layered-substrate corrections are measured
\cite{hansen1988}. It satisfies
\begin{equation}
\label{eq:ch2.scalar-helmholtz-green-pde}
\big(\nabla^{2}+k_0^{2}\big)\,G(\mathbf{r},\mathbf{r}')
=
-\delta(\mathbf{r}-\mathbf{r}'),
\end{equation}
with \(\delta\) the three-dimensional Dirac distribution and boundary conditions on
\(\Gamma\) inherited from \(\mathcal{B}_\Gamma\) \cite{stratton1941}. Equation
\eqref{eq:ch2.scalar-helmholtz-green-pde} is the local symbol of spectral inversion:
for fixed \(\mathbf{r}'\), \(G(\cdot,\mathbf{r}')\) is the response to a unit point
source at \(\mathbf{r}'\) under the same completion as \(\mathcal{L}_\omega\)
\cite{jackson1999}. Convolution with \(G\) therefore propagates an arbitrary scalar source
to the field \(\phi\) that satisfies the inhomogeneous Helmholtz equation on the same domain.
The associated solution operator is
\begin{equation}
\label{eq:ch2.scalar-green-solution-operator}
\phi(\mathbf{r})
=
\int_{\Omega} G(\mathbf{r},\mathbf{r}')\,q(\mathbf{r}')\,d^{3}r',
\qquad
\big(\nabla^{2}+k_0^{2}\big)\phi=q,
\end{equation}
for scalar sources \(q\) in the dual class paired with \(\phi\)
\cite{stratton1941}.
Readers should interpret Eq.~\eqref{eq:ch2.scalar-green-solution-operator} as a weighted
accumulation of responses arriving from every source point inside the domain $\Omega$.
The Green function determines how strongly a source at $\mathbf{r}'$ influences the
observation point $\mathbf{r}$, while $q(\mathbf{r}')$ specifies how much source strength is
available at that location \cite{stratton1941}. In plain language: $\phi(\mathbf{r})$ is the
scalar field one would measure or compute at the observation point; $q(\mathbf{r}')$ is the
source density distributed throughout the impressed region $\Omega$; and $d^{3}r'$ sums
contributions from every infinitesimal source element \cite{jackson1999}. Mathematically,
$\phi(\mathbf{r})$ carries units that depend on the physical quantity represented (potential,
pressure, or a component precursor to $\E$); $q(\mathbf{r}')$ is the corresponding source
density per cubic metre; integration extends only over $\Omega$, the region where sources are
nonzero \cite{stratton1941}.
On exterior domains, $G$ must satisfy outgoing decay compatible
with Section~\ref{sec:ch24}; on bounded domains, Dirichlet, Neumann, or impedance
data on \(\Gamma\) replace the free-space radiation template
\cite{harrington2001}.
Before examining Eq.~\eqref{eq:ch2.free-space-scalar-green} itself, it is useful to recall
the familiar physical picture of a disturbance spreading outward from a localized source.
An expanding spherical wavefront carries phase $e^{-\jj k_0 R}$ and amplitude dilution
proportional to $1/R$; the same geometric spreading responsible for the inverse-square law
in power appears here as the inverse-distance amplitude factor \cite{jackson1999}. Phase
accumulates linearly with propagation distance $R$, and outgoing energy transport requires
the retarded phase sign fixed by the locked phasor convention $\exp(-\jj\omega t)$
\cite{stratton1941}. In unbounded vacuum with outgoing selection at infinity, Sommerfeld
radiation at large $R$ singles out the spherical wave whose amplitude falls as $1/R$ and whose
phase advances with propagation distance.
\begin{equation}
\label{eq:ch2.free-space-scalar-green}
G_0(\mathbf{r},\mathbf{r}')
=
\frac{e^{-\jj k_0 R}}{4\pi R},
\qquad
R=|\mathbf{r}-\mathbf{r}'|,
\end{equation}
with \(R>0\) and the branch of \(k_0\) fixed by the locked phasor axis
\cite{stratton1941}. Equation~\eqref{eq:ch2.free-space-scalar-green}
satisfies Eq.~\eqref{eq:ch2.scalar-helmholtz-green-pde} for
\(\mathbf{r}\neq\mathbf{r}'\); the distributional definition at coincidence is
completed in Subsection~\ref{sec:ch232}. Physically,
Eq.~\eqref{eq:ch2.free-space-scalar-green} is the spherical outgoing wave produced
by a unit monopole source at \(\mathbf{r}'\); its phase gradient defines local
propagation direction on the propagating sheet of Subsection~\ref{sec:ch224}
\cite{stratton1941}. Green kernels realize the resolvent
Eq.~\eqref{eq:ch2.resolvent-operator} developed in Section~\ref{sec:ch22}: the scalar and
dyadic definitions below implement $\mathcal{G}_{E,\omega}$ of
Eq.~\eqref{eq:ch2.helmholtz-solution-operator} and feed the radiation integral
Eq.~\eqref{eq:ch2.radiation-integral} \cite{stratton1941}. Electric-wave operators
and material-weighted symbols from Subsection~\ref{sec:ch222} are presupposed; gauge enters
only through the Lorenz reduction below. Coincidence limits belong to
Subsection~\ref{sec:ch232}; outgoing completion at infinity to Section~\ref{sec:ch24}.
Vector potential and vector kernels. Direct field inversion from measured $\E$ alone is difficult because Maxwell's equations
couple $\E$ and $\H$ through curls, material response, and boundary data; the map from
sources to fields is nonlocal and orientation-dependent \cite{stratton1941}. Vector
potentials are introduced not as additional fields one would place on a probe calibration
sheet but as mathematical intermediaries that simplify reconstruction. In Lorenz gauge,
each Cartesian component of $\mathbf{A}$ satisfies the scalar Helmholtz template of
Eq.~\eqref{eq:ch2.scalar-helmholtz-green-pde} with source $-\mu_0\Jimp$
\cite{jackson1999}. The vector potential is not an
additional physical field but as a mathematical intermediary that simplifies the
reconstruction of electromagnetic quantities; gauge choices remove redundant degrees of
freedom so that propagation integrals take a standard convolution form while what an
engineer measures remains $\E$ and $\H$ \cite{stratton1941}.
Lorenz-gauge vector potentials \(\mathbf{A}\) reduce vector Maxwell inversion to a
vector Helmholtz system driven by \(\Jimp\). Under the gauge and continuity constraints
fixed in Chapter~\ref{ch:01}, each Cartesian component of \(\mathbf{A}\) satisfies
Eq.~\eqref{eq:ch2.scalar-helmholtz-green-pde} with \(k_0\) and source \(-\mu_0\Jimp\)
\cite{stratton1941}.
In Lorenz gauge the vector potential \(\mathbf{A}(\mathbf{r})\) carries units of
T\,m (equivalently Wb\,m\(^{-1}\) in SI); each Cartesian component satisfies the scalar
Helmholtz template with source \(-\mu_0\Jimp\). The electric field phasor follows from
\(\mathbf{A}\) and its divergence through
\begin{equation}
\label{eq:ch2.lorenz-field-from-potential}
\E(\mathbf{r})
=
\jj\omega\mu_0\,\mathbf{A}(\mathbf{r})
+\frac{1}{\jj\omega\varepsilon_0}\,\grad\,\divg\mathbf{A}(\mathbf{r}),
\end{equation}
with \(\E\) in V\,m\(^{-1}\), \(\varepsilon_0\) the vacuum permittivity (F\,m\(^{-1}\)),
and \(\H\) obtained from \(\curl\E\) when the first-order system is required
\cite{harrington2001}. Equation~\eqref{eq:ch2.lorenz-field-from-potential}
is a vector reduction tool, not a substitute for the dyadic definition below when
\(\mathcal{B}_\Gamma\) couples \(\E\) and \(\H\) at boundaries: impedance and PEC
closures are enforced on \(\mathcal{L}_\omega\), not on \(\mathbf{A}\) alone
\cite{stratton1941}. A vector Green function \(\mathbf{G}^{A}(\mathbf{r},\mathbf{r}')\)
may be assembled from three scalar kernels; the electric dyadic is preferred when
\(\Jimp\) is the primary excitation because it maps \(\Jimp\) to \(\E\) directly
without intermediate gauge potentials \cite{harrington2001}.
The scalar Green function describes how a point disturbance propagates. Real electromagnetic
sources, however, possess orientation and polarization. A current element directed along
$\hat{\mathbf{x}}$ generally produces field components in multiple directions. The scalar
kernel is therefore insufficient, and the propagation law must be upgraded to a tensorial
form \cite{stratton1941}. Let \(\overline{\mathbf{G}}_{E}(\mathbf{r},\mathbf{r}')\) denote a rank-two tensor
(dyad) acting on vector sources by dot product from the left; $\mathbf{r}$ is the
observation coordinate and $\mathbf{r}'$ the source coordinate.
The electric dyadic Green function answers, in dyadic form, the impulse-response
question: what electric field \(\E(\mathbf{r})\) is produced at \(\mathbf{r}\) when a
vector current impulse is placed at \(\mathbf{r}'\)?
Method-of-moments assembly implements Eq.~\eqref{eq:ch2.electric-dyadic-convolution}
by testing this kernel against patch basis functions; near-field scanners apply the same
propagation law when translating probe-corrected samples to a reference surface
\cite{harrington2001}. Readers should notice that the dyad encodes polarization
coupling absent from the scalar $G$: a transverse feed excites both radiative and reactive
dyadic components, which is why the Hertzian dipole test at $k_0 a\ll 1$ remains the
lowest-dimensional calibration before aperture catalogs are trusted \cite{jackson1999}.
It inverts the electric-wave operator
Eq.~\eqref{eq:ch2.electric-wave-operator} in the sense
\begin{equation}
\label{eq:ch2.dyadic-green-defining-operator}
\mathcal{W}_{E,\omega}\,
\overline{\mathbf{G}}_{E}(\cdot,\mathbf{r}')
=
\overline{\mathbf{I}}\,\delta(\mathbf{r}-\mathbf{r}'),
\end{equation}
where \(\overline{\mathbf{I}}\) is the identity dyad and the operator
\(\mathcal{W}_{E,\omega}\) acts on the \(\mathbf{r}\)-dependence of
\(\overline{\mathbf{G}}_{E}\) \cite{stratton1941}. Equation
\eqref{eq:ch2.dyadic-green-defining-operator} is the dyadic analogue of
Eq.~\eqref{eq:ch2.scalar-helmholtz-green-pde}: each column of
\(\overline{\mathbf{G}}_{E}\) is the electric field produced by an impulsive vector
source aligned with a Cartesian unit vector at \(\mathbf{r}'\)
\cite{jackson1999}.
Every method-of-moments code and every near-field scanner relies on the convolution
structure below: given an impressed current distribution $\Jimp$, the electric field at
observation points follows by superposition through the dyadic propagator
$\overline{\mathbf{G}}_{E}$. The physical measurement problem is forward prediction---what
field does a known or assumed current produce?---which is the prerequisite for calibration,
coupling analysis, and array synthesis \cite{harrington2001}. Inversion asks the inverse
question from sparse $\E$ samples and is substantially harder even though forward
propagation is linear and explicit \cite{hansen1988}.
Superposition over an extended impressed current \(\Jimp\) then yields the electric field
through the dyadic convolution below; the prefactor \(-\jj\omega\mu_0\) ties the integral
to the locked phasor convention \(\exp(-\jj\omega t)\).
The electric field due to impressed current is
\begin{equation}
\label{eq:ch2.electric-dyadic-convolution}
\E(\mathbf{r})
=
-\jj\omega\mu_0
\int_{\Omega}
\overline{\mathbf{G}}_{E}(\mathbf{r},\mathbf{r}')\cdot\Jimp(\mathbf{r}')\,d^{3}r'
\;+\;\E_{\mathrm{hom}}(\mathbf{r}),
\end{equation}
The importance of Eq.~\eqref{eq:ch2.electric-dyadic-convolution} lies not in the integral
itself but in the asymmetry between prediction and reconstruction. Given $\Jimp$, the field
follows directly. Given only sparse measurements of $\E$, however, recovering $\Jimp$
becomes substantially more difficult \cite{hansen1988}. Here
\(\E_{\mathrm{hom}}\) is the homogeneous solution fixed by \(\mathcal{B}_\Gamma\)
\cite{stratton1941}.
In Eq.~\eqref{eq:ch2.electric-dyadic-convolution}, \(\E(\mathbf{r})\) is the electric
field phasor in V\,m\(^{-1}\), \(\Jimp(\mathbf{r}')\) is the impressed current density in
A\,m\(^{-2}\), \(\mu_0\) is the vacuum permeability (H\,m\(^{-1}\)), and \(\omega\) is the
angular frequency in rad\,s\(^{-1}\); the volume integral extends over the source region
\(\Omega\), and \(\E_{\mathrm{hom}}\) collects excitations not attributable to \(\Jimp\),
including incident fields and cavity modes fixed by \(\mathcal{B}_\Gamma\).
Equation~\eqref{eq:ch2.electric-dyadic-convolution} is the constructive export map every
method-of-moments impedance matrix discretizes: patch currents $\Jimp$ tested against
$\overline{\mathbf{G}}_{E}$ produce the $\E$ that near-field scanners back-propagate to
equivalent surfaces \cite{harrington2001}. The consequence for inversion is
immediate: specifying $\Jimp$ determines $\E$ through a single propagator integral, while
recovering $\Jimp$ from sparse $\E$ samples is ill-posed and requires the regularization
program of Subsection~\ref{sec:ch232} \cite{stratton1941}. Equation~\eqref{eq:ch2.electric-dyadic-convolution} is the
formal realization of Eq.~\eqref{eq:ch2.radiation-integral} with outgoing dyad
\(\boldsymbol{\Gamma}=\overline{\mathbf{G}}_{E}\) on the physical branch
\cite{harrington2001}. The prefactor \(-\jj\omega\mu_0\) matches the locked phasor
convention and the forcing map Eq.~\eqref{eq:ch2.helmholtz-forcing-map}.
The operator $\mathcal{G}_{E,\omega}$ in Eq.~\eqref{eq:ch2.dyadic-green-integral-operator}
should be read as the machine that converts admissible forcing functions into electromagnetic
responses: it accepts a current distribution in the forcing class and returns the
corresponding field throughout the observation domain \cite{harrington2001}. The integral operator \(\mathcal{G}_{E,\omega}\) of
Eq.~\eqref{eq:ch2.helmholtz-solution-operator} is realized by
\begin{equation}
\label{eq:ch2.dyadic-green-integral-operator}
\big(\mathcal{G}_{E,\omega}\mathbf{f}\big)(\mathbf{r})
=
\int_{\Omega}
\overline{\mathbf{G}}_{E}(\mathbf{r},\mathbf{r}')\cdot\mathbf{f}(\mathbf{r}')\,d^{3}r',
\end{equation}
for \(\mathbf{f}\) in the forcing class of Eq.~\eqref{eq:ch2.helmholtz-forcing-map}
\cite{stratton1941}. When \(\mathbf{f}_{E,\omega}\) is the electric Helmholtz drive
induced by \(\Jimp\), Eqs.~\eqref{eq:ch2.electric-dyadic-convolution}
and~\eqref{eq:ch2.dyadic-green-integral-operator} coincide up to the linear map from
\(\Jimp\) to \(\mathbf{f}_{E,\omega}\) \cite{harrington2001}. The solution map
$\mathcal{S}_\omega$ is not merely an abstract mapping notation---it represents the complete
electromagnetic experiment beginning with impressed current and ending with observable
fields, including homogeneous background $\E^{\mathrm{hom}}$ and boundary data
$\mathcal{B}_\Gamma$ when present \cite{stratton1941}. The solution map on
impressed currents is therefore
\begin{equation}
\label{eq:ch2.solution-map-dyadic-realization}
\mathcal{S}_\omega[\Jimp]
=
\mathbf{u}_\omega
\quad\text{with}\quad
\E\ \text{from Eq.~\eqref{eq:ch2.electric-dyadic-convolution}},\
\H\ \text{from}\ \curl\E\ \text{when single-field reduction is valid},
\end{equation}
on domains satisfying Eq.~\eqref{eq:ch2.single-field-reduction-validity}; otherwise
the first-order dyadic system of Section~\ref{sec:ch234} applies
\cite{stratton1941}.
\begin{figure}[t]
\centering
\ChOneFigBlock{%
\begin{tikzpicture}[ch1fig, node distance=5mm and 7mm]
 \node[ch1box] (j) {$\Jimp(\mathbf{r}')$};
 \node[ch1box, right=7mm of j] (g) {$\overline{\mathbf{G}}_{E}$};
 \node[ch1box, right=7mm of g] (e) {$\E(\mathbf{r})$};
 \node[ch1box, below=5mm of g] (bc) {$\mathcal{B}_\Gamma$};
 \node[ch1box, below=5mm of e] (h) {$\H=\mu_0^{-1}\jj\omega^{-1}\curl\E$};
 \draw[ch1flow] (j) -- node[above,font=\scriptsize] {conv.} (g) -- (e);
 \draw[ch1flow] (bc.south) -- (bc.south |- g.north);
 \draw[ch1flow] (bc.south |- g.north) -- (g.north);
 \draw[ch1flow] (e.south) -- (e.south |- h.north);
 \draw[ch1flow] (e.south |- h.north) -- (h.north);
\end{tikzpicture}%
}
\caption{Dyadic inversion chain Eq.~\eqref{eq:ch2.electric-dyadic-convolution}:
impressed current convolves with \(\overline{\mathbf{G}}_{E}\) fixed by
\(\mathcal{B}_\Gamma\); \(\H\) follows from curl relations when reduction is admissible.}
\label{fig:ch2.dyadic-convolution-measurement chain}
\end{figure}
Equation~\eqref{eq:ch2.electric-dyadic-convolution} is a Fredholm integral of the first kind
when \(\mathbf{r}\) is fixed and \(\Jimp\) is unknown, and a constructive export map when
\(\Jimp\) is specified. The prefactor \(-\jj\omega\mu_0\) ties the kernel to the locked
phasor axis \(\exp(-\jj\omega t)\); reversing the sign convention would flip incoming versus
outgoing content without changing the local curl--curl symbol. In vacuum, for \(k_0 R=3\)
with \(R=|\mathbf{r}-\mathbf{r}'|\), the transverse dyadic component of
Eq.~\eqref{eq:ch2.free-space-dyadic-green} contributes \(\approx 70\%\) of
\(|\overline{\mathbf{G}}_{E,0}|\) in operator norm, while the \(R^{-2}\) and \(R^{-3}\)
reactive terms still matter for \(k_0 a\gtrsim 1\) near-zone scans \cite{stratton1941}.
Boundary data \(\mathcal{B}_\Gamma\) enter through the kernel that satisfies the same traces,
not through a separate post-multiplication on \(\E\) alone
\cite{harrington2001}.
In uniform vacuum with outgoing condition at infinity, let
\(\mathbf{R}=\mathbf{r}-\mathbf{r}'\), \(R=|\mathbf{R}|\), and
\(\widehat{\mathbf{R}}=\mathbf{R}/R\).
When observation and source points are separated (\(R>0\)), the outgoing vacuum dyad may
be written explicitly in terms of these geometric quantities. The projector
\(\overline{\mathbf{I}}-\widehat{\mathbf{R}}\widehat{\mathbf{R}}\) represents transverse
radiation: only field components perpendicular to the separation vector carry energy to
infinity in the far zone \cite{stratton1941}. The inverse-power corrections in $k_0 R$
represent stored reactive energy---magnetic and electric energy oscillating locally without
net outward transport at short range. An engineer measuring fields close to an antenna would
observe significant contributions from these reactive terms; far from the source they become
progressively less important, leaving only the transverse radiative component
\cite{harrington2001}.
The standard outgoing electric dyad is
\begin{equation}
\label{eq:ch2.free-space-dyadic-green}
\overline{\mathbf{G}}_{E,0}(\mathbf{r},\mathbf{r}')
=
\frac{e^{-\jj k_0 R}}{4\pi R}
\left[
\overline{\mathbf{I}}
-\widehat{\mathbf{R}}\widehat{\mathbf{R}}
+\left(\frac{1}{\jj k_0 R}-\frac{1}{k_0^{2}R^{2}}\right)
\Big(\overline{\mathbf{I}}-3\,\widehat{\mathbf{R}}\widehat{\mathbf{R}}\Big)
\right],
\end{equation}
valid for \(R>0\) \cite{stratton1941}. The
\(\overline{\mathbf{I}}-\widehat{\mathbf{R}}\widehat{\mathbf{R}}\) term is transverse
to \(\mathbf{R}\) and dominates in the radiative far zone; the $R^{-2}$ and $R^{-3}$ terms
in the parentheses supply near-field reactive structure through inverse powers of \(k_0 R\)
and dominate on electrically small probes placed within a wavelength of the feed
\cite{stratton1941}. Equation~\eqref{eq:ch2.free-space-dyadic-green} satisfies
Eq.~\eqref{eq:ch2.dyadic-green-defining-operator} for
\(\mathbf{r}\neq\mathbf{r}'\); coincidence evaluation is deferred to
Subsection~\ref{sec:ch232}. The same kernel appears in
Eq.~\eqref{eq:ch2.radiation-integral} once \(\mathbf{r}\) lies outside
\(\operatorname{supp}\Jimp\) and outgoing radiation condition is active
\cite{harrington2001}. The magnetic field may be reconstructed from the electric field through Maxwell curl relations
when the observation region is source-free and simply connected: $\H$ is not an independent
measurement channel in the Lorenz reduction but follows from the spatial derivatives of
$\E$ \cite{stratton1941}. This reconstruction is valid in homogeneous regions away from
impressed currents and fails across material interfaces or inside the source support where
curl relations must be supplemented by impressed-current terms \cite{harrington2001}.
The magnetic field may be generated by the electric dyad through Maxwell curl relations
when \(\Jimp\) is divergence-free in the source region:
\begin{equation}
\label{eq:ch2.magnetic-from-electric-dyad}
\H(\mathbf{r})
=
\frac{1}{\jj\omega\mu_0}\,\curl\E(\mathbf{r}),
\qquad
\E\ \text{from Eq.~\eqref{eq:ch2.electric-dyadic-convolution}},
\end{equation}
on simply connected source-free observation regions \cite{stratton1941}. A magnetic
dyadic \(\overline{\mathbf{G}}_{H}(\mathbf{r},\mathbf{r}')\) (A\,m\(^{-1}\)) inverts
Eq.~\eqref{eq:ch2.magnetic-wave-operator} analogously to
Eq.~\eqref{eq:ch2.dyadic-green-defining-operator}, with operator
$\mathcal{W}_{H,\omega}$ in place of $\mathcal{W}_{E,\omega}$: it is the machine that
maps magnetic impressed sources to $\H$, while $\overline{\mathbf{G}}_{H}$ denotes the
corresponding tensorial impulse response \cite{harrington2001}. Duality
\(\varepsilon_0\E\leftrightarrow\mu_0\H\), \(\Jimp\leftrightarrow\mathbf{J}^{\mathrm{(mag)}}_{\mathrm{eff}}\)
maps \(\overline{\mathbf{G}}_{E,0}\leftrightarrow\overline{\mathbf{G}}_{H,0}\) in
vacuum without re-deriving the magnetic-wave symbol. Subsection~\ref{sec:ch231} states
\(\overline{\mathbf{G}}_{H}\) only where magnetic impressed currents are primary;
otherwise Eq.~\eqref{eq:ch2.magnetic-from-electric-dyad} avoids duplicate kernel
tabulations \cite{stratton1941}.
The preceding formulas assume a uniform medium. Real structures seldom satisfy this
assumption. Substrates, radomes, biological tissues, and layered materials alter
propagation and therefore modify the Green response \cite{collin2001}.
In piecewise media, local symbols \(k(\mathbf{r})\), \(\varepsilon(\mathbf{r})\),
\(\mu(\mathbf{r})\) modify Eq.~\eqref{eq:ch2.dyadic-green-defining-operator} to
\begin{equation}
\label{eq:ch2.material-dyadic-defining-operator}
\mathcal{W}_{E,\omega}^{(\varepsilon,\mu)}\,
\overline{\mathbf{G}}_{E}^{(\varepsilon,\mu)}(\cdot,\mathbf{r}')
=
\overline{\mathbf{I}}\,\delta(\mathbf{r}-\mathbf{r}'),
\end{equation}
with \(\mathcal{W}_{E,\omega}^{(\varepsilon,\mu)}\) from
Eq.~\eqref{eq:ch2.material-helmholtz-electric} \cite{stratton1941}. Free-space
Eq.~\eqref{eq:ch2.free-space-dyadic-green} applies patchwise only when
\(k(\mathbf{r})\) is constant in each patch and interface jumps are enforced by
continuity of tangential \(\E,\H\) across shared boundaries
(Eq.~\eqref{eq:ch1.interface-conditions}); multi-patch assembly is
Subsection~\ref{sec:ch234} \cite{harrington2001}. Reciprocity expresses a fundamental
symmetry of passive electromagnetic environments: interchanging transmitter and receiver
does not alter the propagation law itself \cite{stratton1941}. On passive media without
magneto-electric coupling, reciprocity requires that interchanging observation and
source points leaves the dyadic invariant up to transposition:
\begin{equation}
\label{eq:ch2.dyadic-reciprocity-swap}
\overline{\mathbf{G}}_{E}(\mathbf{r},\mathbf{r}')
=
\overline{\mathbf{G}}_{E}^{\mathsf{T}}(\mathbf{r}',\mathbf{r}),
\end{equation}
at fixed $\omega$; the full proof is in Section~\ref{sec:ch28} \cite{stratton1941}.
Equations~\eqref{eq:ch2.electric-dyadic-convolution}--\eqref{eq:ch2.solution-map-dyadic-realization}
hold when \(\Jimp\) lies in the domain of the forcing map, when
\(\mathcal{B}_\Gamma\) is time-independent and identical under phasor reduction, and
when the selected kernel branch matches the resolvent of
Eq.~\eqref{eq:ch2.resolvent-green-equivalence}. They fail for resonant cavities
without quotienting, for incoming rather than outgoing exterior conditions, and for
sources with \(\divg\Jimp\neq 0\) unless charge density \(\rho\) is included in the
complete source tuple of Section~\ref{sec:ch01.1.2} \cite{stratton1941}.
Near \(\mathbf{r}=\mathbf{r}'\), Eqs.~\eqref{eq:ch2.free-space-scalar-green}
and~\eqref{eq:ch2.free-space-dyadic-green} are pointwise undefined; convolution
meaning is distributional (Subsection~\ref{sec:ch232}).
The scalar monopole kernel $G$ supplies the lowest-dimensional propagation law;
the electric dyad $\overline{\mathbf{G}}_{E}$ is the vector generalization demanded
by polarization on real feeds. Subsection~\ref{sec:ch23.signature-dipole} closes the
loop in the opposite pedagogical order: a measured dipole pattern on $\Gamma_R$
precedes the dyadic tabulation, so the laboratory observation remains the anchor
when the full tensor formulas are applied to extended apertures
\cite{jackson1999}. Scalar monopole law $G$, vector dyad
$\overline{\mathbf{G}}_{E}$, and operator realization $\mathcal{G}_{E,\omega}$ are now fixed with
$\mathbf{r}$ as observation and $\mathbf{r}'$ as source throughout; coincidence limits and
retarded causality refine how those kernels may be convolved with extended $\Jimp$
\cite{stratton1941}.
